\section{SECCIÓN 8. HERRAMIENTAS, TÉCNICAS Y METODOLOGÍAS}

\begin{table}[H]
\centering
\caption{Herramientas, Técnicas y Metodologías}
\label{tab:tools}
\begin{tabular}{|p{3cm}|p{4cm}|p{6cm}|}
\hline
\textbf{Herramienta} & \textbf{Uso} & \textbf{Justificación} \\
\hline
SonarQube & Análisis estático del código inicial & Herramienta fiable para análisis estático, detecta code smells y problemas estructurales, seguridad con datos cuantitativos útiles para cuantificar la calidad del sistema \\
\hline
Qlty & Análisis estático y evaluación de deuda técnica & Plataforma en la nube que permite evaluar la calidad del código al realizar análisis estático para detectar malas prácticas, vulnerabilidades y bugs. Se destaca por trabajar con varios lenguajes a la vez, y dar un reporte especializado de deuda técnica. \\
\hline
Jacoco & Medición de coberturas de las pruebas & Permitirá evidenciar la cobertura de las pruebas en Java y obtener métricas cuantitativas de la cobertura de las pruebas en el código generado. \\
\hline
npm audit & Análisis de dependencias & Permitirá revisar si los paquetes/bibliotecas están desactualizadas y revisar dependencias en conflicto dentro del frontend. \\
\hline
Junit & Medir funcionalidad en los endpoints del sistema & Permitirá cuantificar y verificar funcionalidad y métricas de tiempo del sistema. \\
\hline
Jmeter & Medir la carga de estrés en el sistema & Permitirá medir la cantidad de esfuerzo que el sistema ejerce y necesita en su funcionamiento. \\
\hline
OWASP ZAP & Medición de la seguridad del sistema & Permitirá revisar problemas de seguridad en los endpoints y encontrar puntos clave de mejora. \\
\hline
Lighthouse & Medir el performance del frontend & Herramienta fiable para medir la usabilidad del sistema y cuantificar la accesibilidad del mismo, separando y mostrando los puntos fuertes y débiles en el frontend. \\
\hline
Cypress & Medir calidad del frontend & Pruebas de interfaz con métricas reales con pruebas en conjunto al funcionamiento del frontend. \\
\hline
Github & Control sobre el proceso QA & Permitirá revisar y tener un espacio controlado donde revisar el sistema y sus cambios. \\
\hline
\end{tabular}
\end{table}
